\section{Open Challenges and Future Research Directions}
\label{sec:future}

Contemporary legal AI, as the chronological account in Sections~\ref{sec:foundations}--\ref{sec:2026}
has shown, has made genuine progress in legal language understanding, evidence retrieval, document
drafting, and agentic processing. These capabilities, though, remain difficult to bring together
into a system that consistently returns relevant, applicable, and verifiable legal answers. The
central research challenge is therefore not to produce more fluent legal text, but to build systems
that retrieve authoritative evidence, preserve its legal context, reason over relationships between
sources, and verify the claims they generate. Seven directions follow from the synthesis in
Sections~\ref{sec:reliability} and~\ref{sec:jurisdiction}.

\textbf{Knowledge-graph-grounded retrieval.} Flat indexes struggle with questions that need
relationships across documents, such as connecting a provision to its definition and to the
judgments that interpret it. Future work should combine graph-based and text retrieval. It should
keep provenance for every relationship and tell apart explicitly stated facts from inferred
connections. Researchers could evaluate this by comparing a text-only baseline, a graph-only
method, and a combined architecture on multi-hop legal questions.

\textbf{Multilingual and low-resource retrieval.} Researchers should evaluate cross-lingual
retrieval models, terminology alignment, and transliteration-aware indexing in two ways: by
measuring how retrieval performance varies across English, Urdu, and Roman Urdu queries, and by
checking whether legal entity linking reduces errors involving statutory names and section numbers.
The central risk here is mistaking linguistic similarity for legal equivalence.

\textbf{Structure-aware indexing.} Researchers should compare fixed-size, semantic, and hierarchical
chunking under the same conditions. They should pay particular attention to whether preserving
structure improves downstream answer quality beyond its effect on retrieval recall. Parsing
failures also deserve attention: malformed PDFs, inconsistent numbering, and broken cross-references
can all propagate errors through the rest of the pipeline.

\textbf{Claim-level evidence verification.} Citation-grounded generation is not enough if a system
only checks whether a response contains references. Future architectures should tie each material
claim to the specific passage meant to support it. They should resolve citations and flag
unsupported or contradictory statements, because legal support is often more complex than ordinary
textual entailment: a passage may state a general rule while a different passage supplies the
exception.

\textbf{Agentic legal AI with explicit quality gates.} Orchestration should focus on
evidence-controlled workflows rather than simply adding more agents. Each stage needs explicit input
requirements, output contracts, and quality checks. Experiments should compare single-agent and
multi-agent alternatives under equivalent retrieval settings. Without such controls, we cannot
reliably credit improvements to agentic orchestration.

\textbf{Temporal and authority-aware reasoning.} Statutes get amended, and judgments get
distinguished or overruled, over time. Future systems should therefore represent effective dates and
relationships between authorities. They should not assume the newest document is necessarily the
most applicable.

\textbf{Reproducible evaluation and realistic benchmarks.} Future benchmarks should combine
retrieval and generation evaluation while keeping separate scores for each stage, following the
groundedness, retrieval-accuracy, and correctness decomposition this survey described in
Section~\ref{sec:reliability}. They should include difficult cases involving exceptions and
competing authorities. They should also report dataset construction, annotation instructions, and
inter-annotator agreement, so researchers can tell genuine architectural improvements apart from
differences in test construction.

Table~\ref{tab:agenda} pulls these directions together into measurable research objectives.

\begin{table*}[tp]
\centering
\caption{Open challenges and measurable research objectives.}
\label{tab:agenda}
\footnotesize
\setlength{\tabcolsep}{5pt}
\renewcommand{\arraystretch}{1.25}
\begin{tabularx}{\textwidth}{>{\raggedright\arraybackslash}p{3.6cm} >{\raggedright\arraybackslash}X >{\raggedright\arraybackslash}X}
\toprule
\textbf{Research challenge} & \textbf{Proposed direction} & \textbf{Example evaluation} \\
\midrule
Missing cross-document relationships & Combine knowledge graphs with text retrieval. & Recall on multi-hop questions and graph-relation accuracy. \\
Low-resource multilingual queries & Multilingual retrieval with terminology alignment. & Language-specific Recall@$k$ and expert-assessed correctness. \\
Loss of legal hierarchy & Structure-aware indexing and parent-child retrieval. & Structural extraction accuracy and downstream answer quality. \\
Unsupported legal claims & Claim-level evidence verification. & Claim-support precision, recall, and unsupported-claim rate. \\
Unreliable agentic workflows & Modular orchestration with explicit quality gates. & Correctness, failure recovery, latency, and ablation results. \\
Outdated or inapplicable authorities & Temporal and authority-aware retrieval. & Outdated-law error rate and applicability accuracy. \\
Incomparable research results & Shared benchmarks and reproducible protocols. & Cross-system results under common data and evaluation rules. \\
\bottomrule
\end{tabularx}
\end{table*}

These directions depend on one another. Knowledge graphs support authority-aware retrieval only if
source extraction is accurate. Multilingual retrieval improves access but still needs jurisdictional
filtering. Structure-aware indexing preserves context that claim-level verification can then assess.
Agentic workflows can coordinate all of this, but only reproducible evaluation can show whether the
resulting system is actually more dependable. The most promising agenda is therefore an integrated
one: knowledge-graph-grounded, multilingual, structure-aware legal retrieval with explicit evidence
verification and reproducible evaluation. This agenda carries into the next decade the same two
questions this survey opened with in 1991: how should we represent legal knowledge, and how should
we trust a system's reasoning over it?
